\section*{7. Methodik}
Nach einer ausführlichen Literaturrecherche, Prozess- und Anforderungsanalyse, wird das Dashboard entwickelt. Zur Validierung der Eignung als Ablösung des bestehenden Systems wird die Erfüllung der erforderten Funktionalitäten, mittels eines Usability-Tests und einer dazugehörigen Nutzerbefragung anhand beider Systeme, sichergestellt.

Zur Erstellung des Fragenkatalogs wird ein Standard nach ISO 9241-110 auf die gegebenen Umstände angepasst. Hierfür müssen mögliche Antworten und Ergebnisse vorher in Betracht gezogen werden, um die Fragen auf gewünschte Metriken zu reduzieren. Für den Usability-Test werden Testszenarien aufgestellt, welche den Inhalt der Fragen widerspiegeln, damit eine Vollständigkeit der Funktionalität gewährleistet werden kann.

Die Untersuchung erfolgt daraufhin in Form der Auswertung beider Validierungsmaßnahmen. Dadurch können Maßstäbe wie die Effizienz und Nutzerfreundlichkeit des eingeführten Dashboards bewertet werden. Diese Ergebnisse werden schließlich in den umfassenderen Untersuchungskontext eingeordnet, um die Frage der Eignung von Dashboards im Legacy-Kontext zu beantworten.
